\begin{figure}[H]
\centering  
\[
\setlength{\arraycolsep}{2pt}
\begin{array}{cccccccccccccc}

\mathcolorbox{lightorange}{A[1]} 
&\oplus& \mathcolorbox{lightorange}{A[2]} 
&\oplus& \gradientmathbox{lightorange}{lightred}{A[5]} 
&=& \mathcolorbox{lightorange}{100} &\qquad

\gradientmathbox{lightorange}{lightred}{A[5]} 
&\oplus& \gradientmathbox{lightred}{lightpink}{A[6]} 
&\oplus& \gradientmathbox{lightred}{lightpink}{A[8]} 
&=& \mathcolorbox{lightred}{011}\\[1em]

\gradientmathbox{lightred}{lightpink}{A[6]} 
&\oplus& \gradientmathbox{lightred}{lightpink}{A[8]} 
&\oplus& \gradientmathbox{lightpink}{lightpurple}{A[10]} 
&=& \mathcolorbox{lightpink}{110} &\qquad

\gradientmathbox{lightpink}{lightpurple}{A[10]} 
&\oplus& \mathcolorbox{lightpurple}{A[12]} 
&\oplus& \mathcolorbox{lightpurple}{A[15]} 
&=& \mathcolorbox{lightpurple}{111}
\end{array}
\]
\caption{System of linear equations in $\mathbb{F}^3_2$ induced by Figure~\ref{fig:csc592_bigdata_bff}. Each $A[i]$ is the $i$-th entry in the backing array $A$
(zero-indexed). The keys are $\mathcolorbox{lightorange}{x}$, $\mathcolorbox{lightred}{y}$, $\mathcolorbox{lightpink}{w}$ and $\mathcolorbox{lightpurple}{z}$ 
with fingerprints $\mathcolorbox{lightorange}{100}$, $\mathcolorbox{lightred}{011}$, $\mathcolorbox{lightpink}{110}$ and $\mathcolorbox{lightpurple}{111}$, 
respectively. Array entries and fingerprints are marked with corresponding key color(s). There is one equation for each key.}
\label{fig:csc592_bigdata_soe}
\end{figure}